\documentclass[12pt]{article}
\usepackage[utf8]{inputenc}
\usepackage[T1]{fontenc}
\usepackage[margin=1in]{geometry}
\usepackage{setspace}
\usepackage{xcolor}
\usepackage[colorlinks=true, urlcolor=blue!60!black, linkcolor=blue!60!black]{hyperref}
\setlength{\parindent}{2em}
\setlength{\parskip}{1.25em}
\renewcommand{\baselinestretch}{1.0}
\begin{document}
\begin{center}
{\LARGE \textit{Hammer v. Dagenhart}, 247 U.S. 251 (1918)\par}
\vspace{0.8em}
{\large Case Brief --- Four-Page Edition\par}
\end{center}
\vspace{1em}

\section{Memory Jogger}

Congress exceeded its Commerce Clause power and invaded the Tenth Amendment police powers reserved to the States when it barred the interstate shipment of harmless goods produced in factories employing child labor---a dual-federalism holding later expressly overruled in \textit{United States v. Darby} (1941) (\href{https://supreme.justia.com/cases/federal/us/247/251/}{247 U.S. 251}).

\section{Detailed Case Facts}

The Keating-Owen Child Labor Act of 1916 barred interstate or foreign shipment of products of any mine or quarry that had employed children under sixteen within thirty days before removal, and of any mill, cannery, or factory that had employed children under fourteen, or children of fourteen to sixteen for more than eight hours a day, six days a week, or at night. Every state already restricted child labor; North Carolina barred employment of children under twelve.

Roland Dagenhart, a Charlotte cotton-mill worker, had two minor sons working beside him in the mill, one under fourteen and one (Reuben) between fourteen and sixteen. The Act threatened their jobs. On his own behalf and as next friend of his sons, he sued W. C. Hammer, the U.S. Attorney for the Western District of North Carolina, to enjoin enforcement, arguing that the Act was not a valid regulation of commerce and violated the Tenth and Fifth Amendments.

\section{Procedural History}

The Western District of North Carolina held the Act unconstitutional and enjoined its enforcement. Hammer appealed directly to the Supreme Court, which heard argument on April 15--16, 1918, and affirmed on June 3, 1918.

\section{Judicial Votes}

Affirmed, \textbf{5--4}.
\begin{itemize}
\item \textbf{Majority:} Justice Day, joined by Chief Justice White and Justices Van Devanter, Pitney, and McReynolds.
\item \textbf{Concurrences:} None.
\item \textbf{Dissent:} Justice Holmes, joined by Justices McKenna, Brandeis, and Clarke.
\end{itemize}

\section{Holding}

The Keating-Owen Act is unconstitutional. Congress may not use the Commerce Clause to bar interstate transportation of intrinsically harmless goods merely because of the local labor conditions under which they were made. Manufacturing is local, not commerce; Congress may prohibit interstate movement only when the goods are harmful or transportation is necessary to the harm. Using the commerce power to standardize local labor conditions invades the police power reserved by the Tenth Amendment. Day's warnings that ``all freedom of commerce will be at an end'' and his reading of the Tenth Amendment to reserve all powers not ``expressly'' delegated were dicta. A unanimous Court overruled \textit{Hammer} in \textit{Darby}.

\section{Analysis of Opinions}

\textbf{Day (majority).} Drawing on \textit{Gibbons v. Ogden}'s power ``to prescribe the rule by which commerce is to be governed,'' Day distinguished the precedents upholding bans on lottery tickets, adulterated food, prostitution, and liquor: in each, ``the use of interstate transportation was necessary to the accomplishment of harmful results.'' Here the cotton cloth was harmless, and the harm of child labor occurred before shipment. Relying on \textit{Kidd v. Pearson} and \textit{Coe v. Errol}, he drew a categorical line between manufacture and commerce, rejected the Government's theory that Congress could equalize competition among states, and judged the Act by its ``natural and reasonable effect.'' The Court did not reach the Fifth Amendment claim.

\textbf{Holmes (dissenting).} The Act operated only on interstate shipment, squarely within the Commerce Clause, and under \textit{Champion v. Ames} ``regulation means the prohibition of something.'' Citing \textit{McCray}, \textit{Veazie Bank}, and the Sherman Act cases, he insisted that a valid exercise of federal power is not void because of its indirect effects on state policy. The ``harmless goods'' line was arbitrary, and once states send goods across state lines they meet ``the public policy of the United States, which it is for Congress to express.''

\section{Examples: Future Applications}

\begin{enumerate}
\item \textbf{Same-Side (Invalid): Wage-and-Hour Furniture Ban.} Congress bans interstate shipment of furniture made by workers paid below a federal minimum wage. The furniture is harmless and the labor conditions are local, so Congress may not close commerce to equalize labor costs.
\item \textbf{Same-Side (Invalid): Miners' Pensions.} Congress bars interstate shipment of coal from states that do not require pension contributions. Day equated ``the mining of coal'' with making goods, a local activity.
\item \textbf{Opposite-Side (Valid): Toxic Toys and Diseased Meat.} These goods are inherently dangerous and spread harm into the receiving state, within the \textit{Hipolite Egg} exception.
\item \textbf{Opposite-Side (Valid): Stolen Cars.} Interstate transportation is the instrument of the crime, as in \textit{Brooks v. United States} (1925) and the logic of \textit{Champion} and \textit{Hoke}.
\item \textbf{Fence-Sitter: Prison-Made Goods.} Congress bans shipment of convict-made goods only into states that forbid them. The goods are harmless, but like the liquor law in \textit{Clark Distilling} the ban protects receiving states' laws, as later upheld in \textit{Kentucky Whip \& Collar} (1937).
\end{enumerate}

\section{Critique}

\textbf{Originalist and textualist views.} Textualists fault Day for inserting ``expressly'' into the Tenth Amendment, a word the Framers deliberately omitted, as Marshall stressed in \textit{McCulloch}. Yet Randy Barnett and Richard Epstein defend the commerce--manufacture distinction as faithful to Founding-era usage, while Akhil Amar and Jack Balkin read ``commerce'' more broadly and some textualists find Holmes's dissent textually powerful.

\textbf{Progressive and doctrinal views.} From Thomas Reed Powell and Edward Corwin onward, critics condemned \textit{Hammer} as formalism divorced from industrial reality: lottery tickets were not physically dangerous, yet Congress could exclude them on moral grounds. The race to the bottom among states made national regulation the only practical solution, the rationale \textit{Darby} adopted.

\section{Key Quotations}

\begin{quote}
``The making of goods and the mining of coal are not commerce, nor does the fact that these things are to be afterwards shipped or used in interstate commerce make their production a part thereof.''
\end{quote}
\noindent --- Justice Day, Opinion of the Court, 247 U.S. at 272.

\begin{quote}
``In interpreting the Constitution, it must never be forgotten that the Nation is made up of States to which are entrusted the powers of local government. And to them and to the people the powers not expressly delegated to the National Government are reserved.''
\end{quote}
\noindent --- Justice Day, Opinion of the Court, 247 U.S. at 275.

\begin{quote}
``It does not matter whether the supposed evil precedes or follows the transportation. It is enough that, in the opinion of Congress, the transportation encourages the evil.''
\end{quote}
\noindent --- Justice Holmes, dissenting, 247 U.S. at 279--80.

\end{document}
